%========================================
% 12 | Abhiṇha-paccavekkhaṇā-pāṭho
%========================================
\section*{e12 | Abhiṇha-paccavekkhaṇā-pāṭho}\label{sec:e12}

\begin{chantsubsection}
  \chantnote
    {\{Handa mayaṃ abhiṇhapaccavekkhaṇapāṭhaṃ bhaṇāmase\}}
  \chantnotetrans
    {\{Now let us recite the Daily Reflections.\}}

  \chantgap

  \chantnote
    {Jarādhammomhi jarāṃ anatīto\footnote{For gentlemen: \textit{anatīto}; for ladies: \textit{anatītā}.}}
  \chantnote
    {I am of the nature to age; I have not gone beyond aging.}
  \chantindent
    {(rao mi khuam-kae pen tham-ma-da yang-mai luang-phon khuam-kae pai-dai)}

  \chantgap

  \chantnote
    {Byādhidhammomhi byādhiṃ anatīto\footnote{For gentlemen: \textit{anatīto}; for ladies: \textit{anatītā}.}}
  \chantnote
    {I am of the nature to be ill; I have not gone beyond illness.}
  \chantindent
    {(rao mi khuam-jep pen tham-ma-da yang-mai luang-phon khuam-jep pai-dai)}

  \chantgap

  \chantnote
    {Maraṇadhammomhi maraṇaṃ anatīto\footnote{For gentlemen: \textit{anatīto}; for ladies: \textit{anatītā}.}}
  \chantnote
    {I am of the nature to die; I have not gone beyond death.}
  \chantindent
    {(rao mi khuam-tai pen tham-ma-da yang-mai luang-phon khuam-tai pai-dai)}

  \chantgap

  \chantnote
    {Sabbehi me piyehi manāpehi nānābhāvo vinābhāvo}
  \chantnote
    {All that is mine, dear and delightful, will change and be separated from me.}
  \chantindent
    {(rao ja-tong phlat-phrak chak khong-rak khong-chop-jai duai-kan-mod-tang-sin)}

  \chantgap

  \chantnote
    {Kammassakomhi kammadāyādo}
  \chantnote
    {I am the owner of my kamma; the heir to my kamma.}
  \chantindent
    {(rao mi-kam pen khong-tua pen tha-yat haeng-kam)}

  \chantgap

  \chantnote
    {Kammayoni kammabandhu}
  \chantnote
    {Born of my kamma; related to my kamma.}
  \chantindent
    {(mi kam-pen-kam-noet mi-kam-pen-phao-phan)}

  \chantgap
  \clearpage

  \chantnote
    {Kammapaṭisaraṇo\footnote{For gentlemen: \textit{saraṇo}; for ladies: \textit{saraṇā}.}}
  \chantnote
    {I rely on my kamma.\footnote{To rely on kamma: beyond all external supports,
one must return to and bear the results of one's own intentional actions.}}
  \chantindent
    {(mi kam pen-thi-phueng ar-sai)}

  \chantgap

  \chantnote
    {Yaṃ kammaṃ karissāmi}
  \chantnote
    {Whatever kamma I shall do.}
  \chantindent
    {(rao tham-kam-dai-wai)}

  \chantgap

  \chantnote
    {Kalyāṇaṃ vā pāpakaṃ vā}
  \chantnote
    {Whether good or evil,}
  \chantindent
    {(dee-rue-chua ko-tam)}

  \chantgap

  \chantnote
    {Tassa dāyādo\footnote{For gentlemen: \textit{dāyādo}; for ladies: \textit{dāyādā}.} bhavissāmi}
  \chantnote
    {Of that I shall be the heir.}
  \chantindent
    {(rao ja-tong pen-phoo-rap phon-khong-kam-nan)}
    
\end{chantsubsection}
